\documentclass[12pt]{article}
\usepackage[utf8]{inputenc}
\usepackage[T1]{fontenc}
\usepackage{hyperref}
\title{Tela-Masbuar Language Resources for K-12 Curriculum Development}
\author{Compiled by Claude Sonnet 5 (claude-sonnet-5) for Andr\'as Kornai}
\date{}
\begin{document}
\maketitle

\section{Introduction}
Tela-Masbuar (ISO 639-3 tvm; Glottocode tela1241) is an Austronesian language with about 1,000 to 1,100 speakers. It is spoken in the villages of Tela and Masbuar on Babar Island, South Maluku, Indonesia. It is not a Kei or Tanimbar language. My searches found very little for it.

\section{Primary Educational Resources}
None found.

\section{Linguistic and Reference Materials}
\begin{itemize}
\item Encyclopedia article on the language: \url{https://en.wikipedia.org/wiki/Tela%CA%BCa_language}
\item Basic vocabulary wordlist, Austronesian Basic Vocabulary Database: \url{https://lpan.eva.mpg.de/austronesian/language.php?id=315}
\end{itemize}

\section{Religious and Community Resources}
None found.

\section{Digital and Audio Resources}
None found.

\section{Curriculum Development Context}
No teaching materials, orthography or school programme were found. Any curriculum would have to start from documentation, such as fieldwork and a community-led orthography. The wordlist above is the only linguistic data located.
\end{document}